\chapter{Introducción}

Argentina destaca por estar entre los países exportadores agrícolas más importantes del mundo. Según el Banco Mundial, en 2021 el sector agroalimentario
representó aproximadamente el 15,7\% del Producto Interno Bruto (PIB) nacional y generó cerca del 10,6\% de los ingresos fiscales del país \cite{worldbank2024agricultura}. 
La amplia extensión de tierras agrícolas (aproximadamente el 47\% del territorio argentino) sustenta cultivos y producciones que abastecen 
tanto el mercado interno como la exportación \cite{usda2024argentina}. Además, productos como la soja, el maíz y el trigo representan alrededor del 36\% de los ingresos
por exportaciones del país. Al considerar de manera conjunta los productos primarios y las manufacturas de origen agropecuario, el sector alcanza 
aproximadamente el 60\% de las exportaciones totales de Argentina \cite{millermagazine2024argentina}, lo que consolida al sector agrícola como un actor clave para la estabilidad 
económica, la generación de divisas y el sostenimiento del empleo rural.

Las enfermedades en plantas y las plagas constituyen un riesgo permanente   
para los cultivos; en Argentina, los sistemas productivos enfrentan una amplia variedad de enfermedades fúngicas, bacterianas y virales que generan
pérdidas significativas y requieren monitoreo constante para preservar los rendimientos y la calidad de producción \cite{aapresid2024maiz}. La detección tradicional de 
enfermedades se basa en la observación visual en campo y en análisis de laboratorio realizados por organismos como el Servicio Nacional de Sanidad y
Calidad Agroalimentaria (SENASA) \cite{argentina2024laboratorio}.  Estas técnicas permiten identificar síntomas visuales característicos en las hojas, tales como lesiones
necróticas, alteraciones en la coloración y deterioro del tejido vegetal, así como confirmar la presencia de patógenos mediante métodos moleculares,
contribuyendo al manejo integrado de plagas y enfermedades. En paralelo, en los últimos años se han desarrollado herramientas tecnológicas 
orientadas a asistir y automatizar el monitoreo sanitario de los cultivos, incluyendo el uso de sensores, drones y sistemas de análisis de imágenes.
Estas soluciones buscan ampliar la cobertura del monitoreo y reducir la dependencia exclusiva de la inspección manual. Sin embargo, en zonas rurales 
o remotas, donde la conectividad es intermitente o inexistente, el acceso oportuno a plataformas digitales o a sistemas que dependen de infraestructura
externa resulta limitado.

En las últimas décadas, el aumento de la capacidad de cómputo, el acceso a grandes conjuntos de datos y el desarrollo de algoritmos avanzados han 
impulsado la adopción de la Inteligencia Artificial (IA) en una amplia variedad de dominios \cite{dhar2019ondevice}. En el sector agrícola, la creciente demanda de 
agricultura de precisión ha promovido la exploración de tecnologías basadas en IA para mejorar el rendimiento de los cultivos y optimizar el 
monitoreo de enfermedades \cite{ahmed2021mobile,johannes2017automatic,Niaz2025Efficient}. Uno de los enfoques más prometedores es el diagnóstico automatizado de enfermedades foliares, que emplea 
técnicas de visión por computadora, procesamiento de imágenes y algoritmos de aprendizaje automático (Machine Learning, ML) y aprendizaje profundo (Deep Learning, DL) para identificar 
patologías a partir de imágenes \cite{agronomy13040961}. Esta capacidad, sumada a la democratización de la tecnología mediante el uso de dispositivos móviles, 
permite el desarrollo de herramientas accesibles como apoyo en el proceso de inspección visual, sin necesidad de especialistas ni infraestructura 
compleja.

Muchos estudios en los que se implementan modelos de aprendizaje automático y aprendizaje profundo en dispositivos móviles dependen del procesamiento en 
la nube \cite{ahmed2021mobile,Tembhurne2023}, lo que introduce limitaciones como latencia, dependencia de la conectividad y riesgos de privacidad asociados al envío de 
información a servidores externos \cite{xiao2025understanding}. En contraste, la inferencia directa en el dispositivo (on-device), mediante la ejecución del modelo en 
el hardware del teléfono, permite procesar los datos localmente, reduciendo la latencia, preservando la privacidad y habilitando el funcionamiento 
sin conexión \cite{dhar2019ondevice}. Además, al prescindir de recursos en la nube, se disminuyen los costos operativos y se mejora la accesibilidad 
de la solución.

A pesar de estos beneficios, la implementación de modelos de aprendizaje profundo en dispositivos móviles enfrenta retos técnicos significativos.
Los dispositivos móviles tienen recursos limitados de memoria, procesamiento y energía, lo que restringe la complejidad de los modelos que pueden 
ejecutarse sin comprometer la eficiencia. Por eso, es fundamental encontrar un equilibrio adecuado entre precisión del modelo, consumo de recursos
(latencia, energía, memoria) y tamaño de la aplicación.

En este escenario, las arquitecturas ligeras como MobileNet y EfficientNet han demostrado un rendimiento competitivo en tareas de clasificación bajo
restricciones de hardware \cite{zainab2025plant,geetharamani2019identification,subramaniam2026comparative}. Sin embargo, gran parte de los trabajos revisados se concentra principalmente en métricas de precisión del modelo, mientras que son 
más limitados los estudios que analizan de manera comparativa y sistemática el desempeño conjunto de precisión y eficiencia computacional directamente
en dispositivos móviles reales, bajo condiciones controladas y medibles.

Por lo tanto, se vuelve relevante desarrollar un estudio que evalúe comparativamente estas arquitecturas en un entorno de inferencia en el dispositivo, utilizando métricas 
integradas de desempeño (exactitud, sensibilidad y F1-score) junto con métricas de eficiencia (latencia de inferencia, uso de CPU, consumo de memoria, 
energía y tamaño del modelo). Este análisis permitirá determinar la viabilidad técnica real de soluciones basadas en aprendizaje profundo ejecutadas 
localmente en dispositivos móviles, aportando evidencia empírica sobre su aplicabilidad en escenarios rurales con conectividad limitada y 
contribuyendo al diseño de sistemas agrícolas inteligentes con enfoque local.

\section{Objetivos}

\subsection{Objetivo general}

Evaluar la viabilidad técnica del reconocimiento de enfermedades en plantas mediante arquitecturas ligeras de aprendizaje profundo
ejecutadas en dispositivos móviles sin dependencia de conectividad a internet, analizando su desempeño en términos de precisión de
clasificación, eficiencia de ejecución y uso de recursos computacionales.

\subsection{Objetivos específicos}

Para alcanzar este objetivo general, se plantean los siguientes objetivos específicos:

\begin{itemize}
    \item Preparar y preprocesar el conjunto de datos PlantVillage para el entrenamiento, validación y prueba de modelos de clasificación de 
    enfermedades en hojas de plantas.
    \item Implementar y entrenar arquitecturas ligeras de aprendizaje profundo, específicamente MobileNetV3 Small y EfficientNetV2B0, 
    utilizando técnicas de transferencia de aprendizaje y ajuste fino.
    \item Convertir los modelos entrenados a un formato compatible con dispositivos Android mediante TensorFlow Lite, considerando una 
    representación optimizada para su ejecución en el dispositivo.
    \item Desarrollar una aplicación móvil prototipo que permita ejecutar inferencias localmente sobre imágenes de hojas, sin dependencia de 
    servicios externos ni conectividad a internet.
    \item Evaluar comparativamente los modelos en dispositivos móviles con distintas capacidades de hardware, considerando métricas de 
    clasificación, latencia, memoria, consumo energético, temperatura y tamaño del modelo.
    \item Analizar el comportamiento de los modelos frente a imágenes capturadas en condiciones menos controladas que PlantVillage, 
    incluyendo fondos complejos, variaciones de iluminación e imágenes fuera del dominio esperado.
\end{itemize}
